\chapter{Results}
\label{chap:results}

We present results in three stages, following the two-regime story
motivated in Chapter~\ref{chap:method}: first, diagnostic evidence that
appearance/signer identity is measurable in every feature space we test
(\S\ref{sec:results-diagnosis}); second, a low-diversity setting
intervention on NGT744 where synthetic appearance augmentation produces a
clear, sizeable gain (\S\ref{sec:results-lowdiv}); third, a high-diversity setting
intervention on ASL Citizen where the same idea does not lead to improvements(\S\ref{sec:results-highdiv}).
\todome{stale pointer (03.08): ``the two-regime story motivated in
Chapter~\ref{chap:method}'' --- after the agnostic rewrite the Method
chapter no longer carries the regime framing; the motivation now lives in
the Introduction (and the SDDA regime contrast in
\S\ref{sec:related-appearance-bias}). Re-point accordingly. Also missing
here: a space before the final parenthesis.}

\section{Diagnosis: signer identity is recoverable from every feature space}
\label{sec:results-diagnosis}

\subsection{Counterfactual and cross-gloss signer probes (WLASL)}
We first test whether SignCLIP embeddings encode sign content or signer
style, comparing similarity for same-sign/different-signer pairs against
same-signer/different-sign pairs.

\begin{table}[t]
\centering\small
\begin{tabular}{lcc}
\hline
Condition & Shared property & Mean similarity \\
\hline
Positive & Same sign, different signer & \textbf{0.753} \\
Negative & Different sign, same signer & 0.589 \\
Control  & Different sign, different signer & 0.573 \\
\hline
\end{tabular}
\caption[SignCLIP counterfactual analysis on WLASL]{SignCLIP counterfactual analysis on WLASL: mean cosine similarity
by pair type.}
\label{tab:wlasl_counterfactual}
\end{table}

Positive pairs are substantially more similar than control pairs
($\Delta=+0.164$, $p<0.001$), confirming SignCLIP primarily encodes sign
content. Negative pairs (same signer, different sign) are nonetheless
slightly \emph{more} similar than control pairs, indicating a residual
signer-style signal.

We then test signer identifiability directly with a leave-one-gloss-out
retrieval probe: for each query video we build per-signer prototypes
excluding the query's own gloss, and check whether the correct signer is
retrieved.

\begin{table}[t]
\centering\small
\begin{tabular}{lcc}
\hline
Feature & Top-1 acc.\ (chance) & MRR \\
\hline
SignCLIP embeddings & 0.302 (0.013) & 0.471 \\
MediaPipe (1218-d mean+std) & 0.207 (0.013) & 0.374 \\
\hline
\end{tabular}
\caption[Cross-gloss signer retrieval probe]{Cross-gloss signer retrieval probe, SignCLIP embeddings vs.\ raw
MediaPipe pose features.}
\label{tab:cross_gloss_signer_probe}
\end{table}

Both feature spaces identify signers far above chance even with gloss
content excluded, though SignCLIP embeddings do so more strongly. We note
this may partly reflect vocabulary bias (individual signers discussing
different topics) rather than pure appearance signal alone.

\subsection{Signer-separability probes on raw features (ASL Citizen)}
\label{sec:results-l2-raw}

To quantify signer separability in the raw, task-untrained feature spaces,
we apply two complementary probes to mean-pooled per-video features on the
ASL Citizen training split. The \emph{distance probe} computes a per-signer
prototype for every gloss; the intra-signer L2 (video to its own
signer--gloss prototype) measures how consistently a signer repeats a sign,
the inter-signer L2 (between prototypes of different signers for the same
gloss) measures how far apart signers sit at fixed content, and their
per-gloss ratio is summarised by its median over glosses with a $95\%$
bootstrap confidence interval over glosses. The \emph{decoding probe}
trains a linear head to predict the signer from the same mean-pooled
features (3-fold stratified cross-validation; identical protocol to the
fine-tuned-space probes of \S\ref{sec:results-ft-probes}). The two probes
answer different questions: the ratio is a within-space signal-to-noise
statistic (how dominant signer identity is in the untrained geometry),
while decoding measures how much identity information is linearly
accessible. \S\ref{sec:results-ft-probes} shows the two can dissociate
sharply, so we do not read either as a proxy for the other, and in
particular we do not rank spaces by the ratio alone. Because MediaPipe
extraction fails on $8\%$ of training videos, all numbers are computed on
the common subset of $37{,}437$ videos ($33$ signers) for which all three
spaces have features; per-space full-subset values agree within rounding,
so the missing videos drive no comparison.

\begin{table}[t]
\centering\footnotesize
\caption[Signer-separability probes on raw ASL Citizen features]{Signer-separability probes on raw ASL Citizen features (train
split, common subset). Median ratio = per-gloss median of the
inter-signer/intra-signer L2, with the $95\%$ gloss-bootstrap confidence
interval; decoding = linear-head top-1 signer accuracy under 3-fold
stratified cross-validation (chance $3.0\%$, majority-class baseline
$8.0\%$; fold-to-fold sd $\leq 0.5$ points in every space).}
\label{tab:raw_feature_l2_probe}
\begin{tabular}{lrcc}
\toprule
Feature & Dim & Median ratio [$95\%$ CI] & Signer decoding \\
\midrule
MediaPipe        &  609 & $29.9\times$ [28.4, 32.0] & \textbf{98.7\%} \\
Logos (MViTv2-S) &  768 & $22.7\times$ [21.6, 24.0] & 65.9\% \\
I3D (WLASL ft.)  & 1024 & $22.2\times$ [21.2, 23.6] & 94.2\% \\
\bottomrule
\end{tabular}
\end{table}

Signer identity dominates every space by both readouts. The distance
ratio is $22$--$30\times$, with inter $>$ intra for $100\%$ of the
$2{,}731$ glosses, and a linear head recovers the signer far above both
chance ($3.0\%$) and the majority-class baseline ($8.0\%$) in all three
spaces. MediaPipe pose features are the most signer-separable of the
three on \emph{both} probes: their ratio confidence interval overlaps
neither RGB backbone's, and signer decoding reaches $98.7\%$, against
$94.2\%$ for I3D and $65.9\%$ for Logos; the ordering is unchanged on
the signer-disjoint test split ($99.4\%$ / $95.2\%$ / $74.3\%$, $11$
signers, chance $9.1\%$).
\todome{ADDITION START (04.08): gloss-disjoint control, job 25196749}
Nor is the decoding accuracy content leakage: under gloss-disjoint
cross-validation (folds split by gloss, so the classifier is tested
only on signs that appeared for no signer during its training), the
numbers are essentially unchanged ($98.7\%$ / $94.6\%$ / $64.5\%$),
ruling out the possibility that the probe partly learns \emph{which
glosses} a signer contributes rather than how they look and move.
\todome{ADDITION END. Provenance:
probe\_out/gloss\_disjoint\_decoding.json; same features, subset, and
linear head as the table, StratifiedGroupKFold over the 2{,}731 glosses,
zero test videos excluded.}
Camera framing cannot explain this: the pose
features are root-centred at the shoulder midpoint and scale-normalised
to shoulder width at extraction, so what identifies the signer is
shoulder-normalised geometry: body proportions, habitual posture, and
movement style. Pose discards pixels, but it retains, and arguably
distills, precisely the appearance cues that identify signers without
them.

The two probes also dissociate \emph{within} the RGB pair, foreshadowing
\S\ref{sec:results-ft-probes}: Logos and I3D are indistinguishable by
the distance ratio, yet I3D sits near ceiling for linear signer decoding
while Logos gives up $28$ points, consistent with Logos' gloss-supervised,
transfer-oriented pretraining shedding linearly accessible identity
information while leaving the distance geometry untouched. This is an
early instance of a pattern that recurs throughout the thesis: coarse
distance statistics can stay static while the accessible signer
information moves.

\section{Low-diversity regime: NGT744}
\label{sec:results-lowdiv}

\subsection{Which synthetic views help? Sources and dose}
\label{sec:ngt-synthetic-views}
Synthetic views add two things at once: \emph{identities} (avatars or
edits) and \emph{viewpoints} (cameras). These arms separate them by
varying which views form the positive set, their source and their number
$K$. The comparison that isolates appearance is E1 vs.\ E3 vs.\ E4, where
every view sits at the front camera: a head-to-head of the motion-locked
and generative routes at matched $K$, content and viewpoint.

All arms train on the same 744 sentences
(Table~\ref{tab:ngt_setups}), sample one repeat take per step, and hold
the contrastive batch at 64 signs, so in-batch negatives are equal across
arms and only the encoder-pass cost scales with $K$; the learning-rate
schedule follows each arm's update count. Evaluation is fixed: front-view
take-0 galleries and the held-out 364-sentence test set, scored in both
directions with per-sentence ranks so that arms can be compared pairwise
(exact McNemar on top-1 correctness, bootstrap CI on the R@1 difference)
rather than by seed-mean gaps alone.
\begin{table*}[t]
\centering
\footnotesize
\begin{tabular}{llcccc}
\toprule
\textbf{Arm} & \textbf{Training views ($K$ per item)} & \textbf{R@1} & \textbf{R@5} & \textbf{R@10} & \textbf{MRR} \\
\midrule
\multicolumn{6}{l}{\textit{Test $\to$ Train \; (query: test set, gallery: train set)}} \\
\midrule
E1  & 1 real front view (baseline), $K{=}1$ &  0.234\pmsd{0.025} & 0.420\pmsd{0.040} & 0.506\pmsd{0.054} & 0.325\pmsd{0.027} \\
E2  & E1 $+$ 2 real side views (baseline), $K{=}3$ &  0.266\pmsd{0.013} & 0.453\pmsd{0.024} & 0.544\pmsd{0.021} & 0.361\pmsd{0.002} \\
\rowcolor{rowhl}
E3  & E1 $+$ 2 front re-renders, $K{=}3$ &  \textbf{0.649}\pmsd{0.044} & 0.836\pmsd{0.026} & 0.899\pmsd{0.018} & 0.738\pmsd{0.033} \\
E4  & E1 $+$ 2 front edits (signer swap, skin tone), $K{=}3$ &  0.505\pmsd{0.010} & 0.729\pmsd{0.010} & 0.808\pmsd{0.011} & 0.607\pmsd{0.009} \\
E5  & E3 $+$ those 2 re-renders at both side views, $K{=}7$ &  0.437\pmsd{0.018} & 0.676\pmsd{0.019} & 0.783\pmsd{0.032} & 0.553\pmsd{0.018} \\
E6  & E4 $+$ shirt edit, $K{=}4$ &  0.491\pmsd{0.025} & 0.713\pmsd{0.028} & 0.788\pmsd{0.030} & 0.593\pmsd{0.020} \\
E7  & E6 $+$ glasses edit (all 4 edits), $K{=}5$ &  0.461\pmsd{0.023} & 0.691\pmsd{0.015} & 0.771\pmsd{0.023} & 0.568\pmsd{0.012} \\
E8  & E5 $+$ E6's 3 edits, $K{=}10$ &  0.525\pmsd{0.022} & 0.770\pmsd{0.016} & 0.855\pmsd{0.018} & 0.636\pmsd{0.021} \\
E9  & 1 re-render $\times$ 3 views, no real video, $K{=}3$ &  0.296\pmsd{0.022} & 0.545\pmsd{0.014} & 0.662\pmsd{0.027} & 0.414\pmsd{0.018} \\
E10 & E6's 3 edits only, no real video, $K{=}3$ &  0.452\pmsd{0.044} & 0.671\pmsd{0.050} & 0.750\pmsd{0.039} & 0.554\pmsd{0.039} \\

\midrule
\multicolumn{6}{l}{\textit{Train $\to$ Test \; (query: train set, gallery: test set)}} \\
\midrule
E1  & 1 real front view (baseline), $K{=}1$ &  0.247\pmsd{0.035} & 0.446\pmsd{0.043} & 0.547\pmsd{0.037} & 0.347\pmsd{0.037} \\
E2  & E1 $+$ 2 real side views (baseline), $K{=}3$ &  0.420\pmsd{0.016} & 0.721\pmsd{0.018} & 0.805\pmsd{0.005} & 0.551\pmsd{0.014} \\
\rowcolor{rowhl}
E3  & E1 $+$ 2 front re-renders, $K{=}3$ &  \textbf{0.695}\pmsd{0.030} & 0.886\pmsd{0.025} & 0.934\pmsd{0.015} & 0.779\pmsd{0.024} \\
E4  & E1 $+$ 2 front edits (signer swap, skin tone), $K{=}3$ &  0.606\pmsd{0.006} & 0.827\pmsd{0.025} & 0.889\pmsd{0.011} & 0.705\pmsd{0.002} \\
E5  & E3 $+$ those 2 re-renders at both side views, $K{=}7$ &  0.482\pmsd{0.025} & 0.766\pmsd{0.022} & 0.859\pmsd{0.030} & 0.615\pmsd{0.023} \\
E6  & E4 $+$ shirt edit, $K{=}4$ &  0.610\pmsd{0.017} & 0.836\pmsd{0.014} & 0.902\pmsd{0.020} & 0.713\pmsd{0.012} \\
E7  & E6 $+$ glasses edit (all 4 edits), $K{=}5$ &  0.581\pmsd{0.019} & 0.807\pmsd{0.030} & 0.881\pmsd{0.003} & 0.683\pmsd{0.014} \\
E8  & E5 $+$ E6's 3 edits, $K{=}10$ &  0.579\pmsd{0.004} & 0.831\pmsd{0.010} & 0.901\pmsd{0.005} & 0.690\pmsd{0.005} \\
E9  & 1 re-render $\times$ 3 views, no real video, $K{=}3$ &  0.252\pmsd{0.017} & 0.484\pmsd{0.059} & 0.596\pmsd{0.063} & 0.367\pmsd{0.032} \\
E10 & E6's 3 edits only, no real video, $K{=}3$ &  0.513\pmsd{0.034} & 0.755\pmsd{0.012} & 0.837\pmsd{0.020} & 0.625\pmsd{0.020} \\
\bottomrule
\end{tabular}
\caption[Synthetic-view arms on NGT744: cross-signer retrieval]{Synthetic-view arms on NGT744: cross-signer retrieval (0--1
scale), mean$\pm$sd over 3 seeds. Re-renders are motion-locked Unreal
avatars and edits are generative appearance edits, both at the front
camera unless a view is named. E1 has no positive pairs at $K{=}1$ and so
trains with the video-only InfoNCE; every other arm uses SupCon over its
$K$ views.}
% Scale harmonised to 0--1 (2026-07-28) per the todo: all values in this
% table and tab:ngt_signer_count are the summarize_ngt_experiments.py
% percentages divided by 100 (same values, same precision, new
% representation); surrounding prose converted to match. The PNGs
% Axis scales: both ngt_synthetic_views_r1_vs_k.png and ngt_dropoff.png now render
% R@1 on the tables' 0--1 scale (ladder rescaled 2026-07-30).
% A presentation version of the ladder figure (test->train only, descriptive
% point labels, no E-codes needed) lives at thesis/ngt_ladder_presentation.png
% (context/results/ngt_ladder_plot/plot_ngt_ladder_presentation.py) — slides
% only, not included in the thesis.
\label{tab:ngt_synthetic_views}
\end{table*}

\begin{figure*}[t]
\centering
\includegraphics[width=\textwidth]{ngt_synthetic_views_r1_vs_k.png}
\caption[Synthetic-view arms: cross-signer R@1 against views per item]{Synthetic-view arms on NGT744: cross-signer R@1 against the number of
views per training item $K$, by augmentation family (3-seed mean$\pm$sd;
front-view evaluation as in Table~\ref{tab:ngt_synthetic_views}, same 0--1 scale).
Lines trace the
Unreal (E1$\to$E3$\to$E5) and Flux (E1$\to$E4$\to$E6$\to$E7) families from
the shared real-only baseline; isolated points mark the real-multiview
control (E2), the two synthetic-only arms trained without any real video
(E9 Unreal-only and E10 Flux-only, open markers), and the combined arm
(E8). Every arm is drawn at its true $K$, so the five arms that share
$K{=}3$ (E2, E3, E4, E9, E10) sit in one column. The Unreal route peaks
sharply at matched $K{=}3$ and
collapses at full strength ($K{=}7$), whereas the Flux route stays flat as
edits are added.}
\label{fig:ngt_synthetic_views_r1}
\end{figure*}

Figure~\ref{fig:ngt_synthetic_views_r1} plots R@1 against $K$ by
augmentation family, which makes the family structure that
Table~\ref{tab:ngt_synthetic_views} lists row-wise directly visible.
With viewpoint held fixed, both synthetic routes recover a large share of
the cross-signer gap. Over the real-only baseline (E1, test$\to$train R@1
$0.234$), the matched-$K{=}3$ Unreal arm (E3) reaches $0.649$ and the
matched-$K{=}3$ Flux arm (E4) $0.505$: paired gains of $\Delta$R@1
$+0.416$ and $+0.272$ test$\to$train, $+0.448$ and $+0.359$
train$\to$test. Exact McNemar was run on the full-strength arms, where the
substantially smaller gains over E1 (E5 $+0.203$, E6 $+0.257$
test$\to$train) are already significant on $3/3$ seeds in both directions.
At matched strength the motion-locked route wins the head-to-head:
E3 exceeds E4 by $\Delta$R@1 $0.144$ (test$\to$train) / $0.089$ (train$\to$test), significant on
all three seeds test$\to$train and two of three train$\to$test. The generative route therefore
retains roughly $65$--$80\%$ (by
direction) of the Unreal appearance gain at equal $K$, content, and viewpoint:
a substantial but incomplete substitute for the motion-locked renders. Both
synthetic-only arms, trained without any real video (E9 Unreal-only $0.296$, E10
Flux-only $0.452$), at least match the real-only baseline, indicating the avatar
and edited frames carry the discriminative sign content on their own.
Crucially,
the real-multiview control (E2; two \emph{real} extra camera views of the
same signers at the same $K{=}3$) does not reproduce the synthetic gain: it
leaves test$\to$train R@1 near baseline ($0.266$ vs E1 $0.234$) and, while it does lift
train$\to$test ($0.420$ vs $0.247$), stays far below the matched-$K$ synthetic arms in
both directions (E3 $0.649$/$0.695$, E4 $0.505$/$0.606$). Extra viewpoints of the
same identities are therefore not what drives the effect; appearance and
identity diversity is.

Adding synthetic views \emph{beyond} the matched $K{=}3$ does not help, and for
the Unreal route it hurts: the full six-view Unreal arm (E5, $K{=}7$) falls to
R@1 $0.437$, far below its own $K{=}3$ counterpart ($0.649$), whereas the extra
Flux variants (E6--E7) leave the Flux score essentially unchanged. The two
full-strength arms consequently invert the matched-$K$ ordering (E6 $>$ E5,
$\Delta$R@1 $+0.054$/$+0.128$, $2/3$ and $3/3$ seeds), an effect that recurs,
attenuated, on the test-matched subset tier (E6s $>$ E5s, $+0.018$/$+0.065$,
$1/3$ and $2/3$ seeds).

Taken together, the synthetic-view arms separate what viewpoint
diversity contributes from what identity diversity contributes. Adding real
viewpoints does help: the two extra camera views in E2 slightly lift the real-only
baseline (E1) from R@1 $0.234$ to $0.266$ test$\to$train and, far more sharply, from
$0.247$ to $0.420$ train$\to$test ($+0.173$). But the same view budget spent on
identity buys much more: at the identical $K{=}3$, two Unreal avatars
rendered at the unchanged front camera (E3, $0.649$/$0.695$) sit far above the
multiview control in both directions. Nor do synthetic viewpoints stack on
top of synthetic identities: rendering the same avatars from three cameras
(E5, $K{=}7$, $0.437$/$0.482$) lands over $0.2$ R@1 below E3, at
the constant 64-sign negative pool shared by all arms, so negative
starvation cannot account for the drop. The two diversity axes are, however,
not independent: in the signer-count ablation that follows, the
appearance-over-real-viewpoints panel (Table~\ref{tab:ngt_signer_count},
Fig.~\ref{fig:ngt_dropoff}) starts from a stronger real baseline than the
single front view (test$\to$train roughly $0.26$--$0.28$ vs $0.21$--$0.24$, train$\to$test
$0.32$--$0.34$ vs $0.24$--$0.26$ across $S$), and the Unreal gain measured from that
stronger baseline is smaller ($+0.042$ to $+0.109$ R@1 across cells, vs $+0.081$
to $+0.160$ in the appearance $+$ viewpoint panel, E1$\to$E5): real viewpoints already supply part of what the
avatars supply. Real
viewpoints help (mostly train$\to$test), synthetic identities at matched $K$ help
far more, synthetic viewpoints beyond matched $K$ hurt, and viewpoint
diversity partially substitutes for identity diversity.


\subsection{Does the synthetic gain survive real diversity? A fixed-data
real-signer-count ablation}
\label{sec:ngt-signer-count}
% Appearance-only panel provenance: the E3-equivalent arms (ngt_sv_unreal_R1..R4)
% never trained in the first attempt (the jobs died at start on the scratch
% inode quota); retrained 2026-07-29 as Snellius jobs 25042412--25042415
% (3 seeds each) and harvested from ngt_sv_unreal_R*_aggregated.json.
The two-regime reading predicts that the Unreal gain is a property of
\emph{low real-signer diversity}: as genuine signer variety enters
training, synthetic appearance variety should have progressively less to
add. The NGT744 corpus permits a direct test: its training pool comprises
four recording sessions, each a distinct real signer. We propose a
fixed-budget ablation design: we train at
$S \in \{1,2,3,4\}$ real signers, nested (each condition's signers are a
subset of the next), and measure at each $S$ the paired gain
$\Delta(S)$ of two avatar arms over the real-only arm (E1, $K{=}1$): the
appearance-only arm (the synthetic-view E3 setup, $K{=}3$: front camera plus
both avatars' front views) and the full Unreal arm (E5, $K{=}7$).

The load-bearing control is that data volume is fixed while diversity
varies: every condition trains on exactly 180 sentences (the largest even
budget the smallest single-signer pool supports), split evenly across its included signers, with a
single (deterministic, render-complete)
take per sentence, so 180 front-view videos, identical epoch size, and
an identical update budget everywhere, and the take-count differences
between signers (1.5--2.4 takes per sentence) cannot masquerade as a
diversity effect. All arms of a condition use the same sentences and the
same takes; the Unreal arms differ only by attaching avatar views of each
take (the two avatar front views for E3, all six avatar views for E5).


In Table~\ref{tab:ngt_signer_count}, the
real-only arm of each panel barely moves with $S$: spreading the same 180
sentences over one, two, three or four real signers leaves cross-signer
retrieval at R@1 $0.21$--$0.24$ from one real view and $0.26$--$0.28$ from
three (test$\to$train), so real signer variety at constant data volume
does not by itself help. The avatar gain does move, and in the
direction the two-regime account predicts: appearance $+$ viewpoint
shrinks from $+0.145$ at $S{=}1$ to $+0.081$ at $S{=}4$ (test$\to$train),
and appearance over real viewpoints repeats the pattern from a stronger
real baseline with smaller gains throughout. Appearance only stands apart
in size: it is the strongest configuration at every $S$ (R@1
$0.450$--$0.527$ test$\to$train), two to three times the appearance $+$
viewpoint gain, and significant in $3/3$ seeds in every cell.

\begin{table}[t]
\centering
\footnotesize
\begin{tabular}{lccc}
\toprule
$S$ (signers) & real-only arm & $+$avatar arm & $\Delta(S)$ \\
\midrule
\multicolumn{4}{l}{\textbf{Appearance $+$ viewpoint}: 1 real view $\to$ $+$ 6 avatar views \;(E1 $\to$ E5)} \\
\midrule
\multicolumn{4}{l}{\textit{Test $\to$ Train}} \\
1 & 0.237\pmsd{0.007} & 0.382\pmsd{0.022} & $+0.145$ \\
2 & 0.210\pmsd{0.019} & 0.321\pmsd{0.020} & $+0.111$ \\
3 & 0.212\pmsd{0.020} & 0.315\pmsd{0.025} & $+0.103$ \\
4 & 0.224\pmsd{0.011} & 0.305\pmsd{0.029} & $+0.081$ \\
\multicolumn{4}{l}{\textit{Train $\to$ Test}} \\
1 & 0.259\pmsd{0.027} & 0.419\pmsd{0.003} & $+0.160$ \\
2 & 0.236\pmsd{0.017} & 0.392\pmsd{0.058} & $+0.156$ \\
3 & 0.244\pmsd{0.037} & 0.376\pmsd{0.019} & $+0.133$ \\
4 & 0.239\pmsd{0.029} & 0.373\pmsd{0.041} & $+0.134$ \\
\midrule
\multicolumn{4}{l}{\textbf{Appearance only}: 1 real view $\to$ $+$ 2 front avatars \;(E1 $\to$ E3)} \\
\midrule
\multicolumn{4}{l}{\textit{Test $\to$ Train}} \\
1 & 0.237\pmsd{0.007} & 0.527\pmsd{0.044} & $+0.289$ \\
2 & 0.210\pmsd{0.019} & 0.483\pmsd{0.027} & $+0.273$ \\
3 & 0.212\pmsd{0.020} & 0.450\pmsd{0.034} & $+0.238$ \\
4 & 0.224\pmsd{0.011} & 0.482\pmsd{0.061} & $+0.257$ \\
\multicolumn{4}{l}{\textit{Train $\to$ Test}} \\
1 & 0.259\pmsd{0.027} & 0.557\pmsd{0.026} & $+0.298$ \\
2 & 0.236\pmsd{0.017} & 0.524\pmsd{0.040} & $+0.288$ \\
3 & 0.244\pmsd{0.037} & 0.523\pmsd{0.030} & $+0.279$ \\
4 & 0.239\pmsd{0.029} & 0.538\pmsd{0.035} & $+0.299$ \\
\midrule
\multicolumn{4}{l}{\textbf{Appearance over real viewpoints}: 3 real $\to$ $+$6 avatars \;(E2 $\to$ E2$+$av.)} \\
\midrule
\multicolumn{4}{l}{\textit{Test $\to$ Train}} \\
1 & 0.269\pmsd{0.028} & 0.378\pmsd{0.005} & $+0.109$ \\
2 & 0.256\pmsd{0.011} & 0.338\pmsd{0.025} & $+0.082$ \\
3 & 0.257\pmsd{0.032} & 0.337\pmsd{0.035} & $+0.080$ \\
4 & 0.277\pmsd{0.017} & 0.360\pmsd{0.012} & $+0.083$ \\
\multicolumn{4}{l}{\textit{Train $\to$ Test}} \\
1 & 0.322\pmsd{0.040} & 0.429\pmsd{0.015} & $+0.106$ \\
2 & 0.327\pmsd{0.008} & 0.422\pmsd{0.033} & $+0.095$ \\
3 & 0.338\pmsd{0.021} & 0.380\pmsd{0.037} & $+0.042$ \\
4 & 0.344\pmsd{0.014} & 0.394\pmsd{0.021} & $+0.049$ \\
\bottomrule
\end{tabular}
\caption[Real-signer-count ablation on NGT744]{Real-signer-count ablation: cross-signer retrieval R@1 (0--1
scale), mean$\pm$sd over 3 seeds. $\Delta(S)$ is the paired gain of the
$+$avatar arm over the real-only arm in the same panel, from unrounded
means.
Codes name the configurations of
Table~\ref{tab:ngt_synthetic_views} (plus E2$+$av., the E2 arm carrying
the same six avatar views, $K{=}9$), but every arm here is retrained at
the 180-sentence budget, so absolute values differ from that table. Per-seed exact
McNemar ($p<0.05$) holds in $3/3$ seeds for appearance-only everywhere,
$3/3$ for appearance $+$ viewpoint except $S{=}4$ test$\to$train, and
$3/3/2/3$ and $3/2/1/2$ (the two directions) for appearance over real
viewpoints.}
\label{tab:ngt_signer_count}
\end{table}

\begin{figure*}[t]
\centering
\includegraphics[width=.9\textwidth]{ngt_dropoff.png}
\caption[Unreal appearance gain against real-signer count]{Unreal appearance gain $\Delta(S)$ against real-signer count, for the
three panels of Table~\ref{tab:ngt_signer_count}: appearance $+$ viewpoint
(E1 $\to$ E5), appearance only (E1 $\to$ E3), and appearance over real
viewpoints (E2 $\to$ E2$+$avatars), in each retrieval direction. Dashed
lines are least-squares fits, annotated with slope per added signer and
$R^2$. The two panels whose avatars also add viewpoints shrink as real
signer diversity grows, and the gain measured over the three-real-view
baseline is smaller throughout, both consistent with the two-regime
reading that synthetic appearance variety substitutes for absent real
diversity.}
\label{fig:ngt_dropoff}
\end{figure*}

\section{High-diversity regime: ASL Citizen}
\label{sec:results-highdiv}

\subsection{Direct backbone fine-tuning: an informative null}
We begin by fine-tuning the Logos MViTv2-S
backbone itself with cross-entropy and, optionally, an appearance-invariance
term, then re-evaluate frozen features
with dictionary retrieval and a linear probe on the identical gallery.
Training a
contrastive head on top of the frozen features follows in the next
subsection.
Table~\ref{tab:backbone_ft_retrieval} reports the result. Its rows name what
is added to plain gloss cross-entropy: \emph{CE only} is the unchanged
backbone trained with cross-entropy on the real videos under our recipe,
\emph{$+$ synth.\ data} adds the generative appearance edits as extra training items,
\emph{$+$ cosine consistency}, \emph{$+$ CLS-SupCon} and \emph{$+$ SDDA}
each add the named invariance term on top of real and synthetic data,
and \emph{$+$ gloss-SupCon} adds a supervised-contrastive term over
same-gloss clips (\S\ref{sec:method-losses}).

\begin{table*}[t]
\centering
\footnotesize
\setlength{\tabcolsep}{3pt}
\begin{tabular}{lcccc@{\hskip 1.0em}ccc}
\toprule
 & \multicolumn{4}{c}{\textbf{Dictionary retrieval}} & \multicolumn{3}{c}{\textbf{Linear probe}} \\
\cmidrule(lr){2-5}\cmidrule(lr){6-8}
\textbf{Configuration} & \textbf{DCG} & \textbf{R@1} & \textbf{R@5} & \textbf{R@10} & \textbf{R@1} & \textbf{R@5} & \textbf{R@10} \\
\midrule
\multicolumn{8}{l}{\emph{Reference systems}} \\
\;\;ST-GCN (end-to-end)$^{\dagger}$~\cite{aslcitizen}  & 76.37 & 59.52 & 82.68 & 88.13 & --- & --- & --- \\
\;\;I3D (end-to-end)$^{\dagger}$~\cite{aslcitizen}     & 79.13 & 63.10 & 86.09 & 90.86 & --- & --- & --- \\
\;\;SignRep (frozen)$^{\dagger}$~\cite{wongSignRepEnhancingSelfSupervised2025} & 90.84 & 81.37 & 96.11 & --- & --- & --- & --- \\
\;\;Logos (frozen)~\cite{ovodovLogosWellTemperedPretrain2025} & \textbf{92.73} & \textbf{83.59} & \textbf{98.36} & \textbf{99.04} & \tbest{86.33} & \tbest{98.34} & \tbest{99.00} \\
\midrule
\multicolumn{8}{l}{\emph{Logos fine-tuned, top-2 unfreeze}} \\
% Run-id mapping: CE only=run2, +synth=run3, +cosine=run1,
% +CLS-SupCon=run1_clscontrast, +SDDA=run_sdda; top-6: run2_mid,
% run1_mid; full: run_full_ce, run_full_glosscon, run_full_sdda.
% Retest deltas (old generation values in git, experiments.tex
% pre-07-21): retrieval 0.01-0.30 points, probe 0.05-1.84 points; the
% cosine probe value is the least stable (1.84).
\;\;CE only (standard fine-tuning) & 93.58 & 85.57 & 98.49 & 99.08 & \textbf{86.26} & 97.67 & 98.43 \\
\;\;$+$ synth.\ data (aug-as-data) & \textbf{93.59} & \textbf{85.59} & 98.52 & \textbf{99.11} & 85.76 & 97.30 & 98.14 \\
\;\;$+$ cosine consistency & 93.57 & 85.45 & \textbf{98.55} & 99.06 & 82.27 & 95.37 & 96.61 \\
\;\;$+$ CLS-SupCon & 93.08 & 84.43 & 98.34 & 99.03 & 85.10 & 97.26 & 98.08 \\
\;\;$+$ SDDA & 92.66 & 83.45 & 98.32 & 99.03 & 85.88 & \textbf{97.95} & \textbf{98.67} \\
\midrule
\multicolumn{8}{l}{\emph{Logos fine-tuned, top-6 unfreeze}} \\
\;\;CE only (standard fine-tuning) & \tbest{\textbf{93.78}} & \tbest{\textbf{85.95}} & \textbf{98.48} & \textbf{99.09} & \textbf{85.69} & \textbf{97.13} & \textbf{97.77} \\
\;\;$+$ cosine consistency & 93.62 & 85.63 & \textbf{98.48} & 99.05 & 82.45 & 94.54 & 95.76 \\
\midrule
\multicolumn{8}{l}{\emph{Logos fine-tuned, full unfreeze}} \\
\;\;CE only (standard fine-tuning) & 93.37 & 84.96 & 98.49 & 99.10 & 85.29 & 97.55 & 98.44 \\
\;\;$+$ gloss-SupCon & 93.38 & 84.94 & \tbest{\textbf{98.58}} & \tbest{\textbf{99.12}} & 85.38 & \textbf{97.61} & \textbf{98.51} \\
\;\;$+$ SDDA & \textbf{93.39} & \textbf{85.00} & 98.54 & 99.11 & \textbf{85.55} & 97.60 & 98.46 \\
\bottomrule
\end{tabular}
\caption[Frozen-feature retrieval after backbone fine-tuning]{Dictionary
retrieval and linear-probe accuracy on ASL Citizen after backbone
fine-tuning, all on the identical full gallery; row names are explained
in the text. $^{\dagger}$ marks values taken from the cited papers under
the same retrieval protocol, which report retrieval only, and SignRep
does not report R@10. Bold marks the best value per column within a
block, a blue cell the best in that column table-wide. Every cell is a
single training run, so near-ties should be read as equivalent.
}
\label{tab:backbone_ft_retrieval}
\end{table*}

Before the findings, one framing note: the frozen Logos row is not an
achievement of this thesis. It is the unmodified Logos backbone of
Ovodov et al.~\cite{ovodovLogosWellTemperedPretrain2025}, whose authors
never evaluated on ASL Citizen; that it nonetheless exceeds the
benchmark's published end-to-end baselines and SignRep's frozen
representation, with no target-corpus training at all, is a property of
their pretraining. For this thesis it is the baseline: the bar that
every fine-tuning and augmentation arm below attempts to clear.

Four findings stand out. (i) Cross-entropy fine-tuning improves retrieval R@1
from 83.59 (best frozen) to 85.57 at top-2 depth (85.95 at top-6), while
the linear probe is already best \emph{without} any fine-tuning (86.33):
fine-tuning improves retrieval clustering at a small cost to linear
separability. (ii) Depth helps only up to a point: full unfreeze (84.96)
is worse than top-6 (85.95). (iii) No appearance or
signer-invariance objective beats plain cross-entropy at matched depth,
and several are worse: aug-as-data is statistically indistinguishable
from the matched CE run (85.59 vs 85.57, a gap far below the retest
spread quantified below), SDDA below both (83.45, at
frozen-baseline level). The
matched full-unfreeze pair makes this precise
(gloss-SupCon $\approx$ CE only, $\Delta$R@1$<0.1$ on both evals),
and the per-variant consistency ablations replicate it: all four
single-variant runs (glasses/shirt/signer-swap/skin) land in a 0.14-point
band around the aggregate (retrieval R@1 85.48--85.62), so no individual edit
type carries hidden signal. Four of the top-2 configurations were
trained a second time on independently regenerated edits and reproduced
the pattern; retrieval moved by $0.01$--$0.3$ points and probe values by
up to $1.8$, which is the spread the within-block orderings have to be
read against. (iv) Cosine consistency actively damages the
linear probe (probe R@1 drops to 82.27 at top-2 and 82.45 at top-6,
and it is also the least stable value in the table, moving $1.8$ points
between repeats) while its retrieval score
looks ordinary: the positive-only pull distorts the geometry in a way
nearest-neighbour retrieval tolerates but a linear head does not. The
two-stage schedule (\S\ref{sec:setup-training}: the consistency term
switched on only after CE convergence, at a $100\times$ lower learning
rate) fits the same pattern: its retrieval R@1 of 85.02 sits within the
plain-CE range (84.96--85.95) rather than above it, while its probe R@1
of 85.82 avoids the linear-probe damage of single-stage consistency
training. 
Whether the consistency loss in fact pulls augmented copies together is
taken up in \S\ref{sec:results-ft-probes} and
\S\ref{sec:results-aug-geometry}.

\subsection{SignCLIP head on frozen Logos/I3D features}
We next evaluate the appearance edits as training data for a contrastive
SignCLIP-style head trained on top of frozen backbone features
(Table~\ref{tab:asl_citizen_retrieval}).

Rows are named compositionally. \emph{Synth.} marks any training that uses
the Flux appearance edits: \emph{$+$ synth.\ data} adds them to the real
training videos as extra items, \emph{synth.\ data only} trains on the
edits alone, and the pair-SupCon rows use them as explicit positives
(\S\ref{sec:method-objectives}); unmarked rows train on real videos only.
The three pose rows are the published results of the original SignCLIP
model~\cite{jiangSignCLIPConnectingText2024} rather than our runs, which
is why they are reported to two decimals where ours carry three. The
second block searches over head recipes with the backbone held frozen;
the third then carries two of those recipes, the weakest and the best,
onto backbones we had fine-tuned (CE-FT is cross-entropy on the real
videos, aug-FT adds the synthetic edits as data, SDDA-FT adds the SDDA
invariance term, \S\ref{sec:method-losses}). That third block uses a
different frame preprocessing from the rows above, so it carries its own
frozen-backbone rows and is read within the block.

\todome{REWRITTEN 12.08: the previous version claimed ``the head is worth
adding'', which this table cannot show (every row has a head). It now
describes what the recipes differ in, read off the run configs. Two
things surfaced while checking: \emph{from scratch} initialises both
towers from pretrained BERT rather than randomly, and \emph{synth.\ data
only} trains on 450 of the 2{,}731 glosses, which is the real reason it
scores low. A separate check on whether the recipe search used the
preprocessing the Logos authors use is still open.}
The five rows differ in only two respects: how many training stages the
head gets, and which clips it sees. All of them share the architecture,
the frozen features, the contrastive loss, the batch size and the
evaluation. \emph{From scratch} is a single stage with both towers
initialised from pretrained BERT weights, so it is untrained on sign
language rather than randomly initialised. \emph{Fine-tuned} adds a
second stage that resumes the first at a fifth of the learning rate, with
weight decay and a shorter annealing schedule. \emph{$+$ synth.\ data}
changes nothing but the training split, adding $8{,}989$ edited clips to
the $40{,}154$ real ones. \emph{Synth.\ data only} replaces the real
videos with the edits, and is not a like-for-like comparison: the edited
subset covers $450$ of the $2{,}731$ glosses, so that row trains on a
sixth of the vocabulary.

Read this way, the second stage is worth $3.9$ points ($0.750$ to
$0.789$) and the edits roughly one more on top of it ($0.789$ to
$0.800$), while the vocabulary-limited row ($0.760$) says nothing about
how good the edits are. All of these are small beside the choice of
features underneath: the same head reaches $0.32$ on I3D features against
$0.750$ on frozen Logos ones, and the published pose-input SignCLIP
models $0.39$--$0.60$.

Carrying those recipes onto fine-tuned backbones shows that the two do
not stack. A from-scratch head does gain from a fine-tuned backbone ($0.744$
on the matched frozen control against $0.772$--$0.782$), but the best head
recipe does not: it scores highest on the frozen backbone ($0.818$) and
slightly lower on every fine-tuned one ($0.809$--$0.812$). Whatever
backbone fine-tuning has to contribute at this point, the head has
already contributed.

The pair-SupCon rows are the largest effect in the table, and they are a
clean comparison: they resume the same from-scratch head as the
fine-tuned rows, so all of these runs are second stages from one
checkpoint and differ only in what that stage optimises. Keeping the
video--text loss and adding the edits as data carries the head from
$0.750$ to $0.800$. Replacing that loss with SupCon over each video and
its own edits leaves it below where it started: $0.567$ with two edits
and $0.414$ with four. The same clips therefore help or hurt according to
the objective they enter through, and within the SupCon runs more of them
is worse (\S\ref{sec:results-iconicity-main} returns to why).

% TABLE CONVENTION (updated 2026-08-05 per Piotr):
% - \textbf on a data cell = best value in that metric column within its
%   comparison block (ties: all bolded).
% - \tbest{...} (blue cell) = best value in that column across the WHOLE
%   table (per eval type, individual cells; Med.R skipped where all-tied).
% - \rowcolor{rowhl} = the recommended configuration overall (may differ
%   from per-column winners; say so in the caption if it does).
% - Stub column header is "Setup"; setup labels are compositional:
%   {From scratch | Fine-tuned} x {(real only) | + synth. data | , synth.
%   data only}, "synth. data" (not "aug") everywhere.
% Applied so far: T2V table + stub headers of both retrieval tables.
% Remaining tables to be swept for per-block bolding (Phase 2).
\begin{table*}[t]
\centering
\footnotesize
\setlength{\tabcolsep}{3pt}
\begin{tabular}{llcccc}
\toprule
\textbf{Setup} & \textbf{Backbone} & \textbf{R@1} & \textbf{R@5} & \textbf{R@10} & \textbf{Med.\,R} \\
\midrule
\multicolumn{6}{l}{\textit{Published SignCLIP baselines}} \\
SignCLIP from scratch        & MediaPipe pose & 0.39 & 0.71 & 0.84 & 2 \\
SignCLIP fine-tuned          & MediaPipe pose & 0.46 & 0.77 & 0.84 & 2 \\
SignCLIP fine-tuned + linear & MediaPipe pose & \textbf{0.60} & \textbf{0.84} & \textbf{0.89} & 1 \\
\midrule
\multicolumn{6}{l}{\textit{Head-recipe search, backbone frozen}} \\
From scratch          & I3D                  & 0.32  & 0.63  & 0.73  & 3 \\
From scratch          & Logos (frozen) & 0.750 & 0.950 & 0.970 & 1 \\
From scratch $+$ synth.\ data & Logos (frozen) & 0.756 & 0.955 & 0.970 & 1 \\
Fine-tuned            & Logos (frozen) & 0.789 & 0.962 & 0.975 & 1 \\
Fine-tuned, synth.\ data only & Logos (frozen) & 0.760 & 0.957 & 0.971 & 1 \\
Fine-tuned $+$ synth.\ data & Logos (frozen) & \textbf{0.800} & \textbf{0.964} & \textbf{0.976} & 1 \\
\midrule
\multicolumn{6}{l}{\textit{The same recipes on fine-tuned backbones}$^{\ddagger}$} \\
From scratch              & Logos (frozen)   & 0.744 & \textbf{0.966} & \textbf{0.977} & 1 \\
From scratch              & Logos (CE-FT)    & 0.772 & 0.929 & 0.944 & 1 \\
From scratch              & Logos (aug-FT)   & 0.781 & 0.927 & 0.943 & 1 \\
From scratch              & Logos (SDDA-FT)  & \textbf{0.782} & 0.958 & 0.970 & 1 \\
\addlinespace
\rowcolor{rowhl}
Fine-tuned $+$ synth.\ data & Logos (frozen)   & \tbest{\textbf{0.818}} & \tbest{\textbf{0.972}} & \tbest{\textbf{0.982}} & 1 \\
Fine-tuned $+$ synth.\ data & Logos (CE-FT)    & 0.809 & 0.946 & 0.957 & 1 \\
Fine-tuned $+$ synth.\ data & Logos (aug-FT)   & 0.812 & 0.942 & 0.954 & 1 \\
Fine-tuned $+$ synth.\ data & Logos (SDDA-FT)  & 0.810 & 0.964 & 0.974 & 1 \\
\midrule
\multicolumn{6}{l}{\textit{Pair-SupCon head (fine-tune; synth.\ edits as explicit positives)}} \\
$K{=}3$ pair-SupCon (orig.\ $+$ 2 edits)        & Logos (frozen) & 0.567 & 0.887 & 0.934 & 1 \\
$K{=}5$ pair-SupCon (orig.\ $+$ 4 edits)        & Logos (frozen) & 0.414 & 0.767 & 0.858 & 2 \\
\bottomrule
\end{tabular}
\caption[Video-to-text retrieval on ASL Citizen]{Video-to-text
retrieval on ASL Citizen; row naming is explained in the text.
$^{\ddagger}$This block uses a different frame preprocessing from the rows
above, so read it against its own frozen-backbone rows rather than against
the earlier numbers. Bold marks the best value per column within a block,
a blue cell the best in the column, and the shaded row the strongest
configuration.}
\label{tab:asl_citizen_retrieval}
\end{table*}


\paragraph{The effect of augmentations on T2V retrieval performance}
\todome{CORRECTED 12.08: this paragraph said ``fine-tuning alone lifts''
(it is the head's second training stage, not backbone fine-tuning) and
read the synth-only row as evidence that synthetic data adds little,
which it cannot be, since that run sees 450 of the 2{,}731 glosses.}
Table~\ref{tab:asl_citizen_t2v} repeats the comparison in the
text-to-video direction, where precision@$k$ replaces recall@$k$: the
gallery holds several correct videos per gloss query, so recall@1 is
capped by gallery multiplicity however good the model is, while
precision@$k$ measures directly whether the top-ranked videos are
correct. The original SignCLIP work does not report this direction on ASL
Citizen at full scale. The ordering
matches the video-to-text table: the second training stage lifts P@1 from
$0.871$ to $0.895$, adding the edits as data during that stage lifts it
again to $0.919$, the strongest configuration on every metric (Mean~R
$1.128$), and adding them during the first stage instead is worth less
($0.883$). The vocabulary-limited row ($0.884$) is no more comparable
here than it was above. Carried onto the fine-tuned backbones, the same recipe reaches
P@1 $0.925$--$0.935$, mirroring its video-to-text behaviour. Which backbone leads does depend on the direction, however: the frozen
backbone is strongest video-to-text ($0.818$ against $0.809$ on CE-FT)
while CE-FT is strongest text-to-video ($0.935$ against $0.925$), so the
shaded row differs between the two tables.
\todome{DRAFT PROPOSAL END. Interpretive reading (why T2V moves more
than V2T --- the anchor-neighbourhood argument) was cut in the 03.08
cleanup; reconnect it here or leave the paragraph descriptive.}

\begin{table*}[t]
  \centering
  \small
  \makebox[\linewidth][c]{
  \begin{tabular}{llcccc}
  \toprule
  \textbf{Setup} & \textbf{Backbone} & \textbf{P@1} & \textbf{P@5} & \textbf{P@10} & \textbf{Mean R} \\
  \midrule
  From scratch          & Logos (frozen) & 0.871 & 0.854 & 0.786 & 1.228 \\
  From scratch $+$ synth.\ data    & Logos (frozen) & 0.883 & 0.861 & 0.790 & 1.213 \\
  Fine-tuned           & Logos (frozen) & 0.895 & 0.876 & 0.808 & 1.166 \\
  Fine-tuned, synth.\ data only & Logos (frozen) & 0.884 & 0.863 & 0.796 & 1.216 \\
  Fine-tuned $+$ synth.\ data      & Logos (frozen) & \textbf{0.919} & \textbf{0.896} & \textbf{0.829} & \textbf{1.128} \\
  \midrule
  \multicolumn{6}{l}{\textit{The best recipe on fine-tuned backbones}} \\

  Fine-tuned $+$ synth.\ data & Logos (frozen)   & 0.925 & 0.902 & \tbest{\textbf{0.850}} & 1.129 \\
  Fine-tuned $+$ synth.\ data & Logos (CE-FT)    & \tbest{\textbf{0.935}} & 0.909 & 0.844 & 1.106 \\
  Fine-tuned $+$ synth.\ data & Logos (aug-FT)   & 0.930 & \tbest{\textbf{0.910}} & 0.848 & \tbest{\textbf{1.099}} \\
  Fine-tuned $+$ synth.\ data & Logos (SDDA-FT)  & 0.928 & 0.906 & 0.849 & 1.112 \\
  \bottomrule
  \end{tabular}
  }
  \caption[Text-to-video retrieval on ASL Citizen]{Text-to-video
  retrieval on ASL Citizen, same configurations as
  Table~\ref{tab:asl_citizen_retrieval} and the same marking convention.}
  \label{tab:asl_citizen_t2v}
\end{table*}

A qualitative view of the gloss organisation of nine of these spaces,
projected with t-SNE on a shared clip sample, is given in
Appendix~\ref{app:tsne_grid}.


\subsection{Does any objective actually remove signer identity? Distance and
decoding probes on the fine-tuned spaces}
\label{sec:results-ft-probes}
The retrieval numbers above measure task quality, not signer encoding: a
fine-tuned space can match the baseline's accuracy while carrying just as
much identity information. We therefore probe every re-extracted space $z$
with two complementary readouts. The \emph{distance probe} repeats the
construction of \S\ref{sec:results-l2-raw} on each space: per test video the
mean L2 distance to same-gloss clips by the \emph{same} signer
($\bar{d}_{\mathrm{intra}}$) versus by \emph{different} signers
($\bar{d}_{\mathrm{inter}}$), summarised as the median ratio
$\rho=\mathrm{med}\,\big(\bar{d}_{\mathrm{inter}}/\bar{d}_{\mathrm{intra}}\big)$;
$\rho\!\approx\!1$ means signers are no longer geometrically separated at
fixed sign. The \emph{decoding probe} is its supervised complement, using standard
linear-probing methodology~\cite{alainUnderstandingIntermediateLayers2017}: a linear
head $p(s\,|\,z)=\mathrm{softmax}(Wz+b)$ is trained (cross-validated, frozen
$z$) to predict the signer label $s$, reporting top-1 accuracy against the
$1/S$ chance level. The two can dissociate (a distance probe misses
identity information that is linearly decodable but small relative to
gloss variation), so an intervention that genuinely removes appearance
identity must lower both, not merely rescale distances. We run both probes
on the frozen baseline, the top-2 family (CE, aug-as-data,
cosine-consistency plus its per-variant ablations, SDDA), and the SignCLIP
heads, reading each family as a matched block; the depth confound forbids
comparing across them. The backbone-family probes use the test split; the
SignCLIP-head spaces are exported for the train split only, so their probes
run on the matched train split (35 signers) and are reported as a separate
panel (Table~\ref{tab:head_signer_probes}) that must not be compared
row-for-row with the test-split numbers. The full-unfreeze runs are not probed: their
contribution is already settled at the accuracy level
(Table~\ref{tab:backbone_ft_retrieval}: gloss-SupCon adds nothing even at full
capacity), and the compute budget was spent where new evidence was possible,
so the probe conclusions below are scoped to top-2 depth.
\begin{table}[t]
\centering
\small
\begin{tabular}{lcc}
\toprule
\textbf{Space} & \textbf{$\rho$ (inter/intra)} & \textbf{signer decoding} \\
\midrule
Logos (frozen) & 17.67$\times$ & 75.3\% \\
% run-id mapping as in tab:backbone_ft_retrieval
fine-tuned, CE only       & 16.46$\times$ & 64.2\% \\
fine-tuned $+$ synth.\ data & 16.44$\times$ & 62.6\% \\
fine-tuned $+$ cosine consistency & 16.66$\times$ & 56.3\% \\
fine-tuned $+$ SDDA       & 17.12$\times$ & 67.6\% \\
\bottomrule
\end{tabular}
\caption[Signer-identity probes on the fine-tuned spaces]{Signer-identity probes on the re-extracted top-2-family spaces.
$\rho$ = median over multi-signer glosses of the inter-signer/intra-signer
L2 ratio (prototype-based); in every space 100\% of glosses have
inter $>$ intra. Decoding = linear-head top-1 signer accuracy
(3-fold stratified CV; test split, 11 signers, chance 0.091).}
\label{tab:signer_probes}
\end{table}

\begin{table}[t]
\centering
\small
\begin{tabular}{lcc}
\toprule
\textbf{Head space (train split)} & \textbf{$\rho$ (inter/intra)} & \textbf{signer decoding} \\
\midrule
frozen backbone (reference)        & 17.67$\times^{\dagger}$ & 66.9\% \\
From scratch                       & 22.29$\times$ & 22.6\% \\
Fine-tuned                         & 22.35$\times$ & 24.6\% \\
Fine-tuned $+$ synth.\ data        & 22.02$\times$ & 22.4\% \\
$K{=}3$ pair-SupCon                & 23.35$\times$ & 20.1\% \\
$K{=}5$ pair-SupCon                & 26.34$\times$ & 21.2\% \\
\bottomrule
\end{tabular}
\caption[Signer-identity probes on the SignCLIP-head spaces]{The same two probes on the SignCLIP-head spaces. Head features
exist only for the train split, so this panel uses the matched train
split (35 signers, chance 0.029) and is not comparable row-for-row with
the test-split panel above; the frozen-backbone row gives the same-split
decoding reference. $^{\dagger}$Reference $\rho$ from the full-split
computation above. The dissociation sharpens in the heads: they separate
signers \emph{more} geometrically than the backbone spaces
($\rho$ 22.0--26.3$\times$) while making identity far \emph{less}
linearly decodable (20--25\% against 66.9\% on the same
split).\todo{framing check}}
\label{tab:head_signer_probes}
\end{table}


The two probes dissociate exactly as anticipated. The distance probe says
essentially nothing moved: signer separation at fixed sign is enormous in
the frozen space ($\rho = 17.7\times$, inter-signer prototype distances
exceeding within-signer spread for \emph{every} multi-signer gloss) and
remains so after every fine-tuning objective, with the invariance-targeted
runs sitting at $16.4$--$17.1\times$, indistinguishable from plain CE
($16.5\times$) and barely below frozen. The decoding probe tells a
sharper story: a linear head recovers signer identity from the frozen
space at 75.3\% (11-way chance 9.1\%), and every fine-tuned space is
substantially less decodable: CE alone already drops this to 64.2\%,
aug-as-data to 62.6\%, and cosine consistency furthest, to 56.3\%, while
SDDA is the weakest intervention at 67.6\% despite its loss explicitly
targeting signer structure. So identity is not literally frozen in place:
fine-tuning, and consistency training most of all, measurably erodes the
\emph{linearly decodable} component of signer identity, in a direction
this coarse distance statistic is too blunt to register. This reduction
does not translate into better retrieval (Table~\ref{tab:backbone_ft_retrieval}:
consistency training is at or below plain CE on both evals), so removing
decodable signer information and improving retrieval accuracy are
dissociable outcomes here, consistent with the geometry of
\S\ref{sec:results-aug-geometry}, where the same objectives are shown to
pull their own training pairs together without closing the gap on
held-out counterfactuals.

\subsection{Where do the appearance counterfactuals land? Feature-space
geometry of the augmentations}
\label{sec:results-aug-geometry}
\todo{this is too long and too dense. Decide what is the most relevant, degrade the rest to appendix}
The interventions above add little in the high-diversity regime (\S\ref{sec:results-highdiv}),
and two very different geometries could produce a null of this shape. The
synthetic edits might be \emph{redundant} (each augment landing so close
to its source video that it supplies no variation the backbone has not
already absorbed from a 35-signer corpus), or they might carry a
\emph{sim-to-real gap}: occupying a systematically distinct, ``synthetic''
region of the feature space that adds domain shift rather than useful signer
diversity. These are not mutually exclusive; the more interesting possibility
is their conjunction, in which each augment sits near its own original yet all
augments are collectively separable from real video by a small, shared
offset. Closing exactly this kind of real--synthetic gap is the explicit
target of adversarial synthetic-data methods such as
SDDA~\cite{fuSignerDiversitydrivenData2024b}, and pairing an image with its
counterfactual is the basis of counterfactual-contrastive
learning~\cite{roschewitzCounterfactualContrastiveLearning2025}; here we
instead \emph{measure} the geometry directly.

Every video $v$ is embedded by the backbone under analysis into a single
vector $f(v)\in\mathbb{R}^{768}$ (clip features mean-pooled over the video).
Each Flux edit of type
$t\in\{$shirt, glasses, skin~tone, signer~swap$\}$
yields a pair $\big(f(v),f(\tilde{v}_t)\big)$, matched by video id, and the
object of study is the \emph{displacement}
$\delta_t(v)=f(\tilde{v}_t)-f(v)$, the arrow the edit draws in feature
space. We propose four measurements to characterise these arrows; headline
numbers are medians over pairs.

\textbf{(1) Displacement, in units of a real signer change.} A raw cosine
$\cos\big(f(v),f(\tilde{v}_t)\big)$ is not interpretable on its own
(embedding spaces concentrate similarity, and the concentration differs
between spaces), so we calibrate it against a reference built from the
same features and metadata. For each source $v$ we measure its mean cosine
$\bar{c}(v)$ to up to $64$ videos of the \emph{same gloss by different
signers}, the natural magnitude of a genuine signer change for that very
video. We propose the \emph{normalised displacement}
$\mathrm{n.disp}=\big(1-\cos(f(v),f(\tilde{v}_t))\big)/\big(1-\bar{c}(v)\big)$
to measure the fraction of a real signer change that the edit induces. A value of
$0.5$ means the edit moves the embedding half as far as replacing the
person would; $1.0$, exactly as far. A second per-video control,
\emph{different glosses by the same signer}, provides a sign-identity
floor: in every space and on both samples, $\geq 99.9\%$ of augments remain
closer to their source than this floor: no edit moves a video past its
own sign.

\textbf{(2) Sim-to-real separability.} To test for a collectively
``synthetic'' region we train a standard linear classifier (logistic
regression on
standardised features) to label real originals vs.\ augments, and report
ROC-AUC under a protocol we propose for such paired data: $5$-fold
cross-validation \emph{grouped by source video}, so that a
video and its own edits always occupy the same fold. The grouping is what
makes the number meaningful. An original and its augment are
near-duplicates, so if one fell in the training fold and the other in the
test fold, the classifier could recognise the test item purely by its
proximity to a memorised training item, inflating AUC without any
generalising real-vs-synthetic direction; with grouping, AUC is the
probability that an augment of an \emph{unseen} video scores more
``synthetic'' than an unseen real one ($0.5$ = indistinguishable). We
report AUC on raw and on $\ell_2$-normalised features: a gap that survives
normalisation is directional, whereas one that vanishes was mere feature
magnitude.

\textbf{(3) The de-shift test: is the gap a single translation?} Write
$f(\tilde{v}_t)=f(v)+\mu_t+\varepsilon$. If the synthetic region were
nothing but the shared offset $\mu_t$, subtracting it should erase
separability. Within each cross-validation fold we estimate $\hat{\mu}_t$
as the mean displacement of the \emph{training-fold pairs only}, subtract
it from every augment, retrain, and score on the held-out fold;
estimating the shift on the evaluation pairs themselves would trivially
centre them, so the fold-internal protocol keeps the test honest (and
each training pair contributes only one of its two members, chosen at
random, since keeping both would equalise the class means by construction
and pin AUC at $0.5$ mechanically). A collapse of
AUC$_{\text{de-shift}}$ to chance means the entire linearly detectable
synthetic signature was one translation per edit type.

\textbf{(4) Meaning kept, appearance moved?}
Two checks. \emph{gloss@1}: does the augment's nearest real video
(source excluded) show the same sign? If so, the edit preserved meaning; the
video-to-video analogue of the counterfactual check in
\S\ref{sec:results-diagnosis}. \emph{flip}: summarising each signer by the
mean embedding of their videos (prototype), the share of pairs whose
nearest prototype differs between source and augment, meaning the edit visibly
moved identity-related appearance (intended for signer swap; unintended
drift for the local edits).
Because the whole analysis runs per feature space, we repeat it on the
fine-tuned backbones of Table~\ref{tab:backbone_ft_retrieval}
(aug-as-data, cosine-consistency, SDDA, and the CE control) to ask whether
any intervention pulls the augments toward their originals or shrinks the
synthetic region, the geometric companion to the retrieval numbers.

\textbf{Two evaluation samples.} We evaluate every space on two disjoint
pair samples. The first comprises the $6{,}000$ augmented training source
videos, each with all four edits; it maximizes statistical power but is
selection-biased by construction, since augmentation targets were chosen
worst-first: the classes with the fewest training signers and lowest
baseline accuracy (\S\ref{sec:method-diffusion}). The second is a
signer-stratified random sample of $500$ \emph{test-split} videos, augmented
with the identical pipeline; its signers are fully disjoint from the $35$
training signers. For the frozen backbone (which never saw ASL Citizen,
so neither sample raises a leakage concern), the held-out sample serves as
a check that conclusions are not artifacts of the rare-gloss selection. For
the fine-tuned spaces the distinction is substantive: the training pairs are
the very items those models optimized on, so they bound the intervention's
effect from above, while the held-out pairs give the honest generalization
number; and a divergence between the two is itself informative,
separating memorization of specific counterfactual pairs from a learned
appearance invariance. Both sides of every pair share identical
aspect-preserving spatial preprocessing and a matched effective frame rate,
so the measured displacement contains no framing or temporal-sampling
artifact. What we deliberately do \emph{not} remove is the editing
pipeline's own low-level footprint: every augment passes through a resize to
longest-side $512$ and the generator's encode--decode round trip, and so
carries a faint resampling texture even in pixels the edit left untouched.
That footprint is part of the synthetic signal being measured, because any
model trained on these augments consumes exactly the same footprint; we
verified separately that the saved augments preserve the source aspect
ratio.

We expect the four edit types to fall on a clear ordering: shirt and glasses
are local edits that should displace the embedding only slightly, whereas
skin-tone and signer-swap are whole-person edits that should move it much
further, approaching the magnitude of a genuine signer change. The decisive
question for this null is \emph{where} that displacement lands. If
real-versus-synthetic separability is low, the counterfactuals are faithful and the null is one of
saturation: the high-signer ASL Citizen backbone had already learned the
invariance the edits supply, consistent with the marginal
Table~\ref{tab:asl_citizen_retrieval} gain and with the opposite,
low-signer NGT744 result (\S\ref{sec:results-lowdiv}). If instead the
displacement is highly separable, the edits
carry a sim-to-real gap that offsets their benefit in feature space. This is
the mechanistic evidence behind the reading in
\S\ref{sec:discussion-regimes}.
\begin{table}[t]
\centering
\footnotesize
\setlength{\tabcolsep}{4pt}
\begin{tabular}{lccccc}
\toprule
\textbf{Edit type} & \textbf{n.disp} & \textbf{AUC} &
\textbf{AUC$_{\text{de-shift}}$} & \textbf{gloss@1} & \textbf{flip} \\
\midrule
\multicolumn{6}{l}{\emph{Training pairs ($6{,}000$ per edit type, selection-biased)}} \\
\;\;glasses     & 0.46 & 0.993 & 0.50 & 0.58 & 0.67 \\
\;\;shirt       & 0.53 & 0.999 & 0.50 & 0.62 & 0.76 \\
\;\;skin tone   & 0.44 & 0.997 & 0.51 & 0.63 & 0.68 \\
\;\;signer swap & 1.04 & 1.000 & 0.51 & 0.49 & 0.91 \\
\midrule
\multicolumn{6}{l}{\emph{Held-out test pairs ($500$ per edit type, signer-disjoint)}} \\
\;\;glasses     & 0.50 & 0.960 & 0.50 & 0.70 & 0.57 \\
\;\;shirt       & 0.47 & 0.990 & 0.50 & 0.75 & 0.56 \\
\;\;skin tone   & 0.46 & 0.972 & 0.50 & 0.74 & 0.54 \\
\;\;signer swap & 1.05 & 1.000 & 0.50 & 0.59 & 0.77 \\
\bottomrule
\end{tabular}
\caption[Augmentation geometry in the frozen Logos space]{Augmentation geometry in the \textbf{frozen} Logos space
(fine-tuned spaces: Table~\ref{tab:aug_geometry_spaces}). All values are
medians over pairs. n.disp = displacement as a fraction of a real signer
change (measurement 1); AUC = real-vs-synthetic linear separability under
$5$-fold cross-validation grouped by source video ($\ell_2$-normalised
features change no value by more than $0.001$); AUC$_{\text{de-shift}}$ =
the same test after removing the fold-estimated mean displacement
(measurement 3); gloss@1 = share of augments whose nearest real video
(source excluded) shows the same sign; flip = share of augments nearest to
a different signer's prototype than their source.}
% E_signer removed 2026-08-03 per Piotr's decision; anchors that
% falsified the directional reading: chance 0.012, real signer change
% 0.074, edits 0.09-0.14 (probe_out/esigner_anchors.json, job 25195647).
\label{tab:aug_geometry_frozen}
\end{table}

The frozen-space geometry (Table~\ref{tab:aug_geometry_frozen}) settles
the two hypotheses in favour of their \emph{conjunction}, with one twist.
The locality ordering holds: garment and skin edits move an embedding
about half as far as a genuine signer change (normalised displacement
$0.44$--$0.53$), while signer swap moves it exactly as far as one
($1.04$--$1.05$) and flips the nearest signer prototype for $91\%$ of
training pairs: the intended identity change registers. Sign content
survives every edit (gloss retention $0.70$--$0.75$ on the unbiased
held-out pairs; lower, $0.49$--$0.63$, on the training pairs, which by
construction target the hardest classes).
\todome{NUMBERS BUG (03.08, verified against the table above): held-out
gloss@1 is $0.70$/$0.75$/$0.74$ for glasses/shirt/skin tone but
$\mathbf{0.59}$ for signer swap --- the quoted $0.70$--$0.75$ silently
drops the signer-swap value while the sentence claims content ``survives
every edit''. Proposed rewrite: ``Sign content largely survives the
appearance edits (gloss retention $0.70$--$0.75$ on held-out pairs for
glasses/shirt/skin tone) but takes a real hit from the identity change
($0.59$ for signer swap); on the training pairs, which by construction
target the hardest classes, retention is $0.49$--$0.63$.'' The drop is
worth owning: the strongest edit type also costs the most sign content,
which connects to the hallucination limitation.} Yet the same augments are
near-perfectly linearly separable from real video (AUC $0.96$--$1.00$,
unchanged by $\ell_2$-normalisation, so the gap is directional rather
than a magnitude artifact): each counterfactual sits close to its source
while the synthetic population as a whole occupies an identifiable
region.
The twist is how \emph{shallow} that region's signature is.
After subtracting a single mean displacement vector per edit type,
estimated leak-free within each fold, separability collapses to chance
($0.48$--$0.51$) in every space and on both samples: to linear probing,
the entire sim-to-real gap is one shared translation. The individual
arrows do scatter widely around that translation (the shared
component is only $32$--$53\%$ of a typical displacement's length, and
no single principal direction explains more than about $10\%$ of the
displacement cloud's variance), but the scatter carries no class
information; everything a linear probe can use is the mean offset. This
is the mechanistic reading behind the null: the entire
linearly detectable synthetic signature is a single shared offset per
edit type, a direction a classifier can trivially factor out.

\begin{table}[t]
\centering
\footnotesize
\setlength{\tabcolsep}{4pt}
\begin{tabular}{lcccc}
\toprule
\textbf{Space} & \textbf{n.disp} & \textbf{gloss@1} & \textbf{AUC} &
\textbf{AUC$_{\text{de-sh}}$} \\
\midrule
\multicolumn{5}{l}{\emph{Training pairs}} \\
\;\;frozen      & 1.04 & 0.49 & 0.997 & --   \\
\;\;CE          & 1.13 & 0.56 & 0.992 & --   \\
\;\;aug-as-data & 1.04 & 0.91 & 0.991 & --   \\
\;\;consistency & 0.79 & 0.89 & 0.986 & --   \\
\;\;SDDA        & 0.76 & 0.67 & 0.996 & --   \\
\midrule
\multicolumn{5}{l}{\emph{Held-out pairs}} \\
\;\;frozen      & 1.05 & 0.59 & 0.976 & 0.50 \\
\;\;CE          & 1.00 & 0.60 & 0.972 & 0.50 \\
\;\;aug-as-data & 1.17 & 0.62 & 0.949 & 0.48 \\
\;\;consistency & 1.11 & 0.57 & 0.942 & 0.49 \\
\;\;SDDA        & 0.80 & 0.66 & 0.961 & 0.50 \\
\bottomrule
\end{tabular}
\caption[Augmentation geometry across feature spaces]{The same geometry across feature spaces, shown for the
signer-swap edit (the strongest intervention; the local edits show the
same pattern; full grid with uncertainty in
Appendix~\ref{app:aug_geometry_grid}). n.disp and gloss@1
as in Table~\ref{tab:aug_geometry_frozen}; AUC pools all four edit
types; AUC$_{\text{de-sh}}$ is on held-out
pairs only because, unlike n.disp and gloss@1, it shows no
training/held-out asymmetry to report (training values within $0.03$
in every space). Reading: a model that \emph{learned} appearance invariance would
contract n.disp and preserve gloss@1 in \emph{both} panels;
gains confined to the training panel (aug-as-data, consistency) are
memorisation of the specific counterfactual pairs seen in training. The
CE control never trained on augments and is correspondingly symmetric.}
\label{tab:aug_geometry_spaces}
\end{table}


Qualitative companions to these numbers (a t-SNE of the augments around their source videos and the per-edit-type cosine distributions) are shown in Appendix~\ref{app:aug_geometry_grid} (Figs.~\ref{fig:aug_vs_orig_tsne} and~\ref{fig:aug_vs_orig_cosine}).


\subsection{Iconicity analysis}
\label{sec:results-iconicity-main}
The organising
construct of the Logos dataset is the \emph{visually similar sign}
(VSSign): signs that are made with the same hands, movement and location
but mean different things, whether because one form carries several
related meanings or because two separate signs differ too slightly to
tell apart on video. Logos authors show
that \emph{grouping} VSSigns under one label (training the encoder
precisely \emph{not} to separate signs that look alike) yields a
better transfer encoder than keeping them
apart~\cite{ovodovLogosWellTemperedPretrain2025}. A backbone
pre-trained this way is not merely agnostic about form-similar pairs;
collapsing them is part of its supervision, and that metric bias
travels with the frozen features from RSL to our ASL inputs.
To measure that bias we use ASL-LEX~2.0, a linguistic database that codes
each of $2{,}723$ ASL signs on five phonological features (movement,
major and minor location, handshape, sign type) and rates its iconicity,
how much its form resembles its meaning, from $1$ to
$7$~\cite{sehyrASLLEX2021,caselliASLLEX2017}. The idea is close to the
VSSign one, but the use is not: Logos groups look-alike signs in order to
train on them, whereas we only measure, and counting shared features
turns ``looks alike'' into a scale from zero to five instead of a
yes/no. ASL Citizen is also a corpus Logos never trained on, in a
different language, recorded under different conditions, so what the
measurement asks is whether the form grouping survived that transfer.

ASL Citizen contains such form-similar groups too, and they sit at the top
of the iconicity scale. Figure~\ref{fig:iconicity_family} shows one:
GUN1, LIGHTER and REMOTECONTROL agree on all five features (neutral
location, \texttt{a} handshape, no path movement, one-handed) because
each depicts handling the object it names, and all three rate $5.3$--$5.6$
for iconicity. Iconic vocabulary is exactly where distinct meanings
converge on the same form, so how a space treats form-similar pairs
decides what it can say about iconic signs.

\begin{figure}[t]
\centering
\includegraphics[width=\linewidth]{iconicity_family_example.png}
\caption[A high-iconicity form family]{Same hands, same movement, same
location, different words: GUN1, LIGHTER and REMOTECONTROL share all
five ASL-LEX form features and rate $5.3$--$5.6$ of $7$ on deaf-signer
iconicity.}
\label{fig:iconicity_family}
\end{figure}

We quantify this with a \emph{form-dose} measure: bucket every pair of
gloss centroids by how many of the five features the two glosses share,
and trace cosine similarity against that count. Four things follow, and
together they say where form organisation should come from.

\emph{It is built by pretraining, and it travels.} The frozen backbone
responds strongly on its own: mean centroid similarity climbs from
$0.01$ at zero shared features to $0.18$ at five (Spearman $\rho=.06$).
Logos never trained on ASL Citizen, so this structure has crossed to
another language, another set of signers and another recording setup. It
is also what the Logos authors credit for the quality of their encoder,
and on this benchmark the untouched frozen backbone still beats every
published system (Table~\ref{tab:backbone_ft_retrieval}).

\emph{In-domain training moves it, and the objective decides which way.}
Fine-tuning the backbone on gloss labels flattens the response, plain
cross-entropy most of all ($0.09$ at five shared features, $\rho=.00$):
the model reorganises around whatever separates this particular
vocabulary and gives up the phonological structure it started with. SDDA
is the one intervention that meaningfully protects it, keeping about half
the frozen slope ($0.13$, $\rho=.03$), while aug-as-data and the
consistency pull sit level with plain cross-entropy.
Figure~\ref{fig:form_dose_rho} shows the pattern across every space we
trained: the SignCLIP heads stay near the frozen level on frozen
features, and inherit the flattening when built on a fine-tuned backbone.

\emph{Pushing appearance invariance hard drives it back up, without
being asked to.} The pair-SupCon heads are the only spaces above the
frozen level ($0.25$ at $K{=}3$, $\rho=.09$; $0.39$ at $K{=}5$,
$\rho=.16$), and they see no gloss labels at all: their positives are the
$K$ views of one source video, and every other clip is a negative, two
GUN1 recordings included (\S\ref{sec:method-objectives}). Nothing tells
them which clips share a word. Collapsing appearance variation simply
removes the axis the space was mostly using, and form is what is left to
organise it, so form-similar glosses drift together as a side effect.

\emph{Overshooting the pretrained level is what costs performance.}
Retrieval falls as $K$ rises, from $0.800$ V2T R@1 for the same edits
used as ordinary training data to $0.567$ at $K{=}3$ and $0.414$ at
$K{=}5$: the space collapses the very contrasts the task must tell apart.
The opposite error is harder to see here, because plain cross-entropy
flattens the response furthest and still scores best on this benchmark
($+2.4$ R@1 over frozen), so on retrieval alone the loss looks free.
Retrieval is not the only readout of those features, though. A linear
classifier on the frozen features reaches $86.33$ R@1, the best probe
score in the table, and no fine-tuned space improves on it. The frozen
geometry is therefore not arranged for nearest-neighbour retrieval, which
is why fine-tuning lifts that score, but the class structure is already
there and linearly separable: fine-tuning rearranges the space to suit
the retrieval metric more than it adds information.
A second reason comes from the Logos paper itself. Its authors
report that co-training on all three of their corpora transfers better to
a low-resource target than fine-tuning their RSL-only model on that same
target ($66.82$ against $65.57$ on WLASL), and account for it only by the
higher accuracy reached~\cite{ovodovLogosWellTemperedPretrain2025}.
Flattening offers a mechanism they never test: fine-tuning on one target
gives up the form structure the pretraining built, whereas co-training
keeps the source objective, and with it that structure, alive.
\S\ref{sec:discussion-regimes} develops this. It stays a conjecture,
because the intermediate checkpoints needed to test it are not public.

\begin{figure*}[t]
\centering
\includegraphics[width=\textwidth]{form_dose_main.png}
\caption[Form-dose response across training recipes]{Centroid cosine
similarity against the number of shared ASL-LEX form features. Boxes are
quartiles, the dashed line the bucket mean. The frozen backbone sits between two opposite
pressures: fine-tuning the backbone flattens the response, most
completely with plain cross-entropy, while training a head with the
appearance counterfactuals as explicit positives amplifies it beyond
the frozen level and increasingly so with $K$. The pair-SupCon panels
also start higher at zero shared features, so the claim concerns the
slope rather than the absolute level. Mean centroid similarity at five
shared features: $0.18$ frozen, $0.09$ CE, $0.13$ SDDA, $0.25$ at
$K{=}3$, $0.39$ at $K{=}5$; the six buckets hold $546{,}736$ to
$2{,}178$ centroid pairs, from zero to five shared features.}
\label{fig:form_dose_ft}
\end{figure*}

\begin{figure}[t]
\centering
\includegraphics[width=\linewidth]{form_dose_rho.png}
\caption[Form organisation across every trained space]{Form organisation
of every space we trained, as the Spearman correlation between shared
form features and centroid cosine over all gloss pairs. The dashed line
is the frozen backbone. Every backbone fine-tune falls below it, heads
trained on frozen features stay near it, heads on fine-tuned backbones
inherit the flattening, and only the pair-SupCon heads, the ones trained
with the appearance counterfactuals as explicit positives, rise above
it.}
\label{fig:form_dose_rho}
\end{figure}

